\documentclass[11pt]{article}

\usepackage[margin=1in]{geometry}
\usepackage{booktabs}
\usepackage{array}
\usepackage{longtable}
\usepackage{hyperref}

\title{The ResearchClassification Repository: A Cross-Walk Between\\
ANZSRC (FOR/SEO), OpenAlex, CWTS Leiden, ASJC and the SDGs}
\author{}
\date{}

\begin{document}
\maketitle

% NOTE TO AUTHOR: insert formal citations (Australian Bureau of Statistics ANZSRC
% documentation; OpenAlex topic classification documentation; CWTS Leiden Ranking
% methodology papers; Elsevier ASJC documentation; UN SDG indicator framework) at the
% points marked [CITE] below, using this paper's own bibliography style.

\section{Introduction}

Bibliometric analyses that combine publication-level data from OpenAlex with
administrative research-output data coded under the Australian and New Zealand Standard
Research Classification (ANZSRC) face an immediate practical obstacle: the two
classification systems were built independently, at different levels of granularity, by
different institutions, for different purposes. OpenAlex organises the scholarly
literature into a four-level hierarchy of domains, fields, subfields and topics, derived
computationally from citation and text data. ANZSRC's Fields of Research (FOR) and
Socio-Economic Objective (SEO) classifications organise Australian and New Zealand
research funding, higher-degree completions and other administrative collections into
divisions, groups and fields (FOR) or divisions, groups and objectives (SEO), revised
periodically by expert working groups (1998, 2008 and 2020 vintages). Neither system was
designed with the other in mind, and no official concordance between them exists.

This document describes a software repository, \texttt{ResearchClassification}, built to
solve this problem systematically: to provide a single, auditable, versioned \emph{cross-walk}
between every major research classification scheme a bibliometric or research-policy
analyst is likely to encounter --- ANZSRC FOR (1998/2008/2020), ANZSRC SEO
(1998/2008/2020), OpenAlex (domain/field/subfield/topic), the CWTS Leiden Ranking's
five-way field classification, Elsevier's All Science Journal Classification (ASJC), and
the United Nations Sustainable Development Goals (SDGs). The repository exposes this
cross-walk through a single Python API, \texttt{Resolver.resolve(code, from\_scheme,
to\_scheme)}, backed by a set of git-tracked, human-readable CSV ``bridge'' tables and a
pre-built DuckDB database for fast lookup. Every correspondence the repository asserts
carries an explicit confidence score and a \texttt{match\_method} tag describing exactly
how it was derived --- hand-curated fact, exact lexical match, algorithmic lexical
cascade, or (where the underlying data genuinely does not support a definitive answer) a
documented, flagged uncertainty. The design goal throughout has been to make every
correspondence traceable to its evidence, rather than to present a single opaque lookup
table.

This document has three purposes: to describe the raw data sources the repository draws
on (Section~\ref{sec:data}); to describe the procedures used to build the correspondences
between them, including a substantial piece of recent work extending the cross-walk from
OpenAlex to FOR2020 down to the finest possible level of granularity on both sides
(Sections~\ref{sec:architecture}--\ref{sec:extension}); and to report the final coverage
achieved and discuss the residual, honestly-documented limitations that remain
(Sections~\ref{sec:limitations}--\ref{sec:results}), together with their implications for
combining OpenAlex-derived data with FOR/SEO-coded administrative data
(Section~\ref{sec:discussion}).

\section{Classification Schemes and Raw Data Sources}
\label{sec:data}

\subsection{ANZSRC FOR and SEO}

The Australian and New Zealand Standard Research Classification (ANZSRC) is jointly
maintained by the Australian Bureau of Statistics and Statistics New Zealand [CITE]. Two
of its component classifications are used here: Fields of Research (FOR), organising
research by subject matter, and Socio-Economic Objective (SEO), organising research by
the purpose it serves. Both have been revised twice: FOR1998/SEO1998 (a flat six-digit
code space, coarser levels zero-padded, e.g.\ division \texttt{210000}, group
\texttt{230100}, leaf \texttt{230101}), FOR2008/SEO2008 (natively variable-width:
division \texttt{01}, group \texttt{0101}, field \texttt{010101}), and FOR2020/SEO2020,
the current, also variable-width, vintage and the repository's canonical target: FOR2020
comprises 23 divisions, 213 groups and 1{,}967 fields (2{,}203 codes); SEO2020 comprises
19 divisions, 128 groups and 840 objectives (987 codes). Official ABS correspondence
tables map each legacy vintage's leaf-level codes forward to FOR2020/SEO2020, but ---
critically for joint use with OpenAlex --- none extends beyond leaf level (no official
division- or group-level correspondence exists), and no official correspondence exists
at all between any FOR/SEO vintage and any non-ANZSRC scheme.

\subsection{OpenAlex}

OpenAlex classifies every work into a four-level hierarchy: domain (4 categories), field
(26), subfield (252) and topic (4{,}516), each nesting exactly inside its parent, derived
computationally from citation-network structure and natural-language processing over
titles, abstracts and citation context. Each topic carries, beyond its label, a short
list of representative keywords and a one-paragraph summary. The repository's raw source
for this hierarchy is a single spreadsheet, \texttt{OpenAlex\_topic\_mapping\_table.xlsx}
(4{,}516 rows, one per topic, with topic/subfield/field/domain identifiers and labels,
keywords, summary and a Wikipedia URL also used to link into the Leiden classification).
It is tracked directly in the repository's data directory precisely because the keyword
and summary text is not carried through into the derived, lighter-weight per-level CSV
files, yet remains needed by the fine-grained matching described in
Section~\ref{sec:extension}.

\subsection{CWTS Leiden Ranking classification}

The CWTS Leiden Ranking assigns every publication to one of five broad ``main fields''
(Social Sciences and Humanities, Biomedical and Health Sciences, Physical Sciences and
Engineering, Life and Earth Sciences, and Mathematics and Computer Science) via an
intermediate layer of algorithmically-derived ``micro-clusters,'' each linked, where
possible, to a representative Wikipedia article. The repository links Leiden
micro-clusters to OpenAlex topics via an exact join on that shared Wikipedia URL
(successful for 3{,}880 of 4{,}516 topics; the remainder have no unique link on one or
both sides and are left unjoined rather than forced into a spurious match), and derives
Leiden main-field correspondences for FOR2020 divisions and groups by majority vote over
that join.

\subsection{ASJC and the Sustainable Development Goals}

Elsevier's All Science Journal Classification (ASJC, 361 codes) is linked to OpenAlex
fields via an exact identifier join (278 correspondences), not a lexical procedure,
since OpenAlex's own field identifiers were originally derived from ASJC. The UN's 17
Sustainable Development Goals and their five ``P'' pillars (People, Planet, Prosperity,
Peace, Partnership) are linked to SEO2020 via a user-supplied, hand-reviewed
division-level alignment table, since SEO --- unlike FOR --- is an objective-of-research
classification with a natural conceptual relationship to development objectives.

\section{System Architecture}
\label{sec:architecture}

Every correspondence, regardless of how it was derived, is expressed in one common
eleven-column schema (a ``bridge'' table): \texttt{source\_system}, \texttt{source\_code},
\texttt{source\_label}, \texttt{system} (target scheme), \texttt{canonical\_code},
\texttt{canonical\_label}, \texttt{canonical\_level}, \texttt{is\_primary},
\texttt{match\_method}, \texttt{confidence} and \texttt{notes}. A source code may have
several ranked candidate targets recorded, with exactly one flagged \texttt{is\_primary}
as the answer \texttt{resolve()} returns by default. \texttt{match\_method} is drawn from a
fixed, documented vocabulary (\texttt{exact\_match}, \texttt{contains\_match},
\texttt{constrained\_lexical}, \texttt{below\_floor}, \texttt{manual\_curated},
\texttt{manual\_override}, \texttt{derived\_empirical}, \texttt{cultural\_proxy},
\texttt{identity}, and a few others), so consumers can filter or re-weight results by
provenance rather than treat every correspondence as equally authoritative.

These bridge tables and each scheme's canonical code table are stored as git-tracked CSV
files --- every change to every correspondence is visible in version control as a diff
--- and loaded into a single pre-built DuckDB database bundled with the Python package.
\texttt{Resolver()} opens this database read-only (\textasciitilde10\,ms), falling back
automatically to rebuilding an in-memory database from the CSVs if the pre-built file
cannot be opened (e.g.\ a DuckDB version mismatch). \texttt{resolve(code, from\_scheme,
to\_scheme)} always requires both schemes named explicitly, never inferred: code values
collide across schemes and vintages in ways that make silent inference unsafe (FOR1998's
\texttt{300101} denotes ``Soil Physics'' while FOR2020's own \texttt{300101} denotes an
unrelated ``Agricultural biotechnology diagnostics''; roughly 48\% of FOR1998's 898 codes
collide this way). \texttt{resolve()} returns a \texttt{CanonicalResult} recording the
resolved code, label, level, match method, confidence and ranked alternates, or raises
\texttt{LookupError} with an explanatory message rather than returning nothing.

\section{Correspondence-Building Methodology}
\label{sec:methodology}

Two broadly different methodologies are used across the repository, and their asymmetric
application to the FOR2020$\leftrightarrow$OpenAlex relationship specifically is the
starting point for the extension described in Section~\ref{sec:extension}.

\subsection{Hand-curated correspondences}

Where a definitive, reviewable correspondence could be established directly --- the
FOR2020 division/group to OpenAlex domain/field/subfield direction, ported from an
earlier, independently hand-curated project; the SEO2020-to-SDG division alignment,
supplied and confirmed by the user; the OpenAlex-to-ASJC identifier join --- the
repository uses that source directly, tagging it \texttt{manual\_curated} or
\texttt{user\_provided} with confidence 1.0 (or, for the ASJC/OpenAlex identifier join,
\texttt{exact\_match}). No further inference is applied.

\subsection{Algorithmic correspondences: the lexical cascade}

Where no hand-curated source exists, correspondences are built algorithmically using a
shared matching module, \texttt{cascade\_match.py}. Given a source label and a pool of
candidate target labels, the cascade tries three strategies in order, stopping at the
first decisive result: (1) exact match between the two labels' stop-word-filtered,
spelling-normalised word \emph{sets} (not a similarity score --- ``Engineering''
\{engineering\} must not match ``Aerospace engineering'' \{aerospace, engineering\}, a
narrower concept); (2) ``contains'' match, where one label's word set is a unique proper
subset of another's (catching compound labels, e.g.\ OpenAlex's ``Pediatrics,
Perinatology and Child Health'' against ANZSRC's ``Paediatrics''); and (3), last resort,
the raw \emph{count} of words shared between the source's descendant-label bag and each
coarse-level candidate's own bag --- a raw count, not a normalised ratio, since earlier
ratio-based formulas (overlap coefficient, Jaccard) were found to systematically favour
whichever candidate had the smaller vocabulary, sometimes egregiously (FOR2020 division
45, Indigenous Studies, by far the largest bag of any division, would otherwise win most
ratio-based comparisons on size alone regardless of relevance). A minimum raw-overlap
floor (\texttt{MIN\_OVERLAP}), below which even the top-scoring candidate is untrusted,
is calibrated separately per granularity: coarser, larger candidate pools use a higher
floor than finer, smaller ones, since a single shared word is far more likely to be
genuine signal in a small pool.

Where the cascade cannot resolve an item with confidence, it is recorded with
\texttt{match\_method = "below\_floor"} and confidence 0.0 --- never silently dropped, and
never forced to a low-confidence guess presented as reliable. Every correspondence,
confident or not, is written to the bridge table, so downstream consumers choose their
own confidence threshold rather than inheriting the pipeline's.

Division 45 (Indigenous Studies) merits a specific note, since it recurs throughout the
repository's history. FOR2020 introduced Indigenous Studies as a first-class division
for the first time (earlier vintages scattered a handful of Indigenous- and
M\=aori-related codes across otherwise-unrelated divisions); its labels are almost
entirely a generic FOR2020 concept with a population prefix added (``Aboriginal and
Torres Strait Islander history,'' a M\=aori-language label with an English gloss,
``Pacific Peoples archaeology''). With no direct OpenAlex, Leiden or ASJC counterpart,
and an aggregate vocabulary large enough to distort any bag-based matching, it is
resolved via a dedicated \texttt{cultural\_proxy} mechanism that routes each division-45
code through the non-Indigenous FOR2020 concept it is a population-specific instance of,
tagged distinctly so the result is never mistaken for an official ANZSRC relationship,
and is excluded from every general-purpose candidate pool used elsewhere.

\section{Extending the OpenAlex$\rightarrow$FOR2020 Direction to Full Precision}
\label{sec:extension}

\subsection{Motivation and the pre-existing asymmetry}

Prior to this work, the repository's FOR2020$\rightarrow$OpenAlex direction
(division/group to domain/field/subfield) was complete and hand-curated, but the reverse
direction --- OpenAlex to FOR2020, the direction actually needed to classify an
OpenAlex-indexed publication into a FOR code for comparison against
administratively-coded research output --- was materially shallower. An OpenAlex
subfield or topic input capped out at FOR2020 \emph{group} (4-digit) precision; an
OpenAlex field input capped out at \emph{division} (2-digit) precision; and no
correspondence existed at all from an OpenAlex topic to a FOR2020 \emph{field} (6-digit
leaf), the finest level on both sides and the level most useful for comparing
OpenAlex-derived and FOR-coded data beyond the broadest disciplinary level.

Two further asymmetries were discovered while doing this work. First, the existing
subfield-to-group bridge, while nominally covering all 252 subfields, was of markedly
uneven quality: 46 subfields (18\%) had no confident match (\texttt{below\_floor}), and
among the remaining 206 the median confidence was only 0.22 --- suggesting a substantial
further share of the ``confident'' picks were also wrong. Second, no field-level
definitional text exists anywhere in the ANZSRC source data for FOR2020 (only division-
and group-level definitions were ever published), so matching against FOR2020 fields has
only the field's own short label to work with, while OpenAlex offers considerably richer
text (topic label, keywords, prose summary) --- a structural asymmetry designed around
rather than treated as an oversight.

\subsection{Step 1: auditing OpenAlex field to FOR2020 division (26 codes)}

The coarsest of the three new correspondences, OpenAlex field to FOR2020 division, was
already complete before this work began, but --- given its foundational role for
everything finer --- was given a full, direct review of all 26 rows rather than taken on
trust. Reading each field's own child subfields against its assigned division, rather
than relying on the label-level lexical score alone, surfaced three corrections:
``Economics, Econometrics and Finance'' had been assigned to division 35 (Commerce,
Management, Tourism and Services) despite FOR2020 possessing a dedicated division 38
(Economics) whose own group ``Econometrics'' exactly matches part of the field's own
label; ``Energy'' (whose child subfields are unambiguously engineering/power-technology
concepts) had been assigned to division 51 (Physical Sciences) purely because the bare
word ``energy'' happened to be a subset of the physics group label ``Particle and high
energy physics,'' an unrelated concept and a textbook contains-match false positive; and
``Decision Sciences'' was redirected, at the user's explicit direction, from division 46
(Information and Computing Sciences) to division 35, reflecting the judgement that this
genuinely cross-disciplinary field (spanning statistics, operations research,
information systems and management science) is best understood as a management
discipline for this cross-walk's purposes.

\subsection{Step 2: a full audit of OpenAlex subfield to FOR2020 group (252 codes)}

Given the quality issues identified above, all 252 OpenAlex subfields --- not merely the
46 originally unmatched --- were reviewed directly, using a purpose-built diagnostic
tool (\texttt{audit\_oax\_for\_bridges.py}) that, for each subfield, displays its current
match, its ranked lexical alternates, and the full list of candidate FOR2020 groups
under both its own ``home'' division and, where the current pick crosses into a
different division, that division's groups too. Reviewing subfields in this format ---
rather than a further round of formula-tuning, an approach the project's own development
history had already found unproductive for a closely analogous earlier problem ---
allowed each assignment to be confirmed, corrected via an explicit, individually
justified override, or (had any remained genuinely undecidable) escalated for a manual
tie-break. In the event, every one of the 252 subfields was resolved through direct
review; none required escalation, well inside the pre-agreed threshold of fewer than 20
unresolved cases.

This review added 47 new manual overrides (on top of approximately 44 already present),
correcting both the 46 originally-unmatched subfields (e.g.\ ``Horticulture,'' corrected
to the exact-fit group ``Horticultural production'') and nine further subfields whose
nominally-confident match was a clear lexical false positive analogous to the ``Energy''
case above (e.g.\ ``Nuclear and High Energy Physics'' had been matched to ``Astronomical
sciences'' rather than the obviously appropriate ``Particle and high energy physics'').
Following this review, zero subfields remain \texttt{below\_floor}: all 252 resolve to a
real, confident FOR2020 group, with a match-method composition of 82 exact matches, 79
constrained-lexical matches and 91 manual overrides (median confidence 0.7).

\subsection{Step 3: clustering OpenAlex topics down to FOR2020 fields (4,516 codes)}

With the subfield-to-group correspondence now solid, the finest-grained correspondence
--- OpenAlex topic to FOR2020 field --- became tractable by a ``clustering down''
strategy: for a given topic, the candidate pool of FOR2020 fields is restricted to only
those fields belonging to the FOR2020 group already matched to that topic's own
subfield (an average of 9.2 candidates, ranging from 1 to 32, rather than the full
1{,}967-field space). This constraint is what makes the finest-grained matching problem
tractable at all, and depends entirely on the subfield-to-group correspondence
underlying it being trustworthy --- precisely why that correspondence was audited to
completeness first.

Each topic is matched against its constrained pool using its own label together with its
OpenAlex-supplied keywords and summary text (sourced by identifier join against the raw
source spreadsheet, avoiding a join on label text, which would be vulnerable to a small
number of known character-encoding corruptions present in roughly 1--2\% of the raw
file's topic labels), scored against each candidate field's own label with the same
cascade described in Section~\ref{sec:methodology}, adapted to normalise by the
candidate field's own short word count rather than the topic's, since the field label is
consistently the thinner side of the comparison at this granularity. Groups with only a
single field are assigned directly, with no scoring required.

Per an explicit joint decision with the user, no fixed manual-review threshold
(analogous to the fewer-than-20-cases rule at the subfield level) was applied at the
topic level, given its far larger scale; genuinely uncertain matches are instead flagged
via \texttt{below\_floor} status and low confidence directly in the data, to be handled
programmatically (Section~\ref{sec:limitations}) rather than reviewed individually. Of
the 4{,}516 topics, 3{,}918 (86.8\%) resolve to a real, confident FOR2020 field; 598
(13.2\%) do not clear the matching floor and are recorded as \texttt{below\_floor}. Two
individual corrections were made following a targeted spot-check of a small sample,
discussed in Section~\ref{sec:limitations}.

\subsection{Integrating the new tiers into \texttt{resolve()}}

The public \texttt{resolve()} method was restructured to try the finest OpenAlex
precision the input supports first, cascading to progressively coarser tiers ---
topic$\rightarrow$field, subfield$\rightarrow$group, field$\rightarrow$division ---
stopping at the first tier not itself flagged \texttt{below\_floor}. A
\texttt{below\_floor} result at any tier is therefore never returned directly; it remains
visible in that tier's own bridge CSV, but the public API always prefers a trustworthy
coarser answer over an untrustworthy finer one. Because the subfield-to-group tier now
has zero \texttt{below\_floor} results, this coarser fallback is itself always reliable:
no OpenAlex topic input can ever fall all the way through to division-level output, a
property verified exhaustively by the automated test suite (Section~\ref{sec:testing}).

This restructuring also surfaced and corrected two latent defects in the previous,
two-tier implementation: the match method returned to the caller had been hard-coded to
a fixed string at each tier rather than read from the underlying bridge row, so every
OpenAlex-to-FOR2020 result previously reported a provenance tag that did not necessarily
reflect how it was actually derived; and the group-level fallback check had been written
so that it could never trigger for a \texttt{below\_floor} result (every subfield already
had a primary row on record, \texttt{below\_floor} or not), meaning the documented
graceful-degradation behaviour had, in practice, never executed for the group tier prior
to this work.

\subsection{Validation}
\label{sec:testing}

Two consistency checks were added alongside the new correspondences. A hard,
exhaustively-checked invariant confirms that every topic's resolved FOR2020 field sits,
without exception, inside the FOR2020 group already assigned to that topic's own
subfield --- guaranteed by the clustering-down design itself, and treated as a permanent
regression guard rather than a data-quality report: any violation could only indicate an
implementation error. A second, deliberately non-fatal check reports how often a
subfield's matched group belongs to a different division than its own parent field's
division; this ``crossing'' is expected and legitimate (OpenAlex's ``Education''
subfield sits under field ``Social Sciences,'' mapping to division 44, yet correctly
belongs with division 39's own Education groups) and occurred, after the audit, for 62
of the 252 subfields (down from approximately 143 before the audit, consistent with the
audit correcting spurious cross-division assignments rather than merely filling gaps).

\texttt{tests/test\_resolver.py}, the repository's plain-assertion test suite, was
extended with six new tests: exhaustive coverage of all 26 OpenAlex fields, 252
subfields and 4{,}516 topics against FOR2020, including the explicit assertions that
zero subfield-level results are ever \texttt{below\_floor} and that no topic-level result
ever falls as far as division level; the topic-nests-under-subfield structural
invariant; and a targeted regression test for the match-method hard-coding defect. All
29 tests in the suite pass against the final, rebuilt database.

\section{Known Limitations and Residual Uncertainty}
\label{sec:limitations}

Two distinct kinds of residual error were identified in a small, informal spot-check of
the 598 \texttt{below\_floor} topics conducted after the main extension was complete,
and are recorded here explicitly rather than corrected, since correcting them was judged
out of scope for the present work and the distinction between them is itself
methodologically informative for anyone extending or auditing this cross-walk further.

The first kind is a genuine \emph{within-pool scoring miss}: the topic ``Surface
Roughness and Optical Measurements'' sits under the OpenAlex subfield ``Computational
Mechanics,'' whose matched FOR2020 group, ``Fluid mechanics and thermal engineering,''
does in fact contain a well-fitting field, ``Non-Newtonian fluid flows (incl.\
rheology)'' --- the correct answer is present in the candidate pool the topic-clustering
procedure constructed, but the lexical scoring procedure failed to select it. This kind
of error is, in principle, fully recoverable by refining the scoring procedure alone
(or, more practically, by a targeted large-language-model judgement call restricted to
the small candidate pool already identified, discussed further in
Section~\ref{sec:discussion}), without requiring any change to the upstream
subfield-to-group correspondence.

The second kind is structurally different and cannot be fixed by better scoring at all:
the topic ``Biodiesel Production and Applications'' sits under the OpenAlex subfield
``Biomedical Engineering,'' whose matched FOR2020 group, ``Biomedical engineering,''
contains no biofuel- or biotechnology-related field whatsoever among its candidates
(biomaterials, biomedical imaging, tissue engineering, medical devices, and similar ---
none a plausible home for a topic about biodiesel production). Here, the mismatch lies
in OpenAlex's own topic-to-subfield placement, upstream of anything the present
cross-walk pipeline controls; no amount of re-scoring within the constrained candidate
pool can produce a correct answer, because a correct answer does not exist within that
pool. Both failure modes present identically in the bridge CSV --- both are tagged
\texttt{below\_floor} with confidence 0.0, and both cause \texttt{resolve()} to fall back
identically to the subfield's group-level answer --- so distinguishing them requires
manually tracing a given topic back through its own subfield to that subfield's
candidate-field pool, exactly the diagnostic context the repository's audit tooling
already surfaces for the subfield-to-group tier but does not yet surface, as a
standalone tool, for the topic tier.

Because this distinction was identified from a sample of only ten \texttt{below\_floor}
topics (of which two, or 20\%, exhibited one or other of these failure modes), it should
be treated as suggestive rather than as a reliable estimate of the true residual error
rate across all 598 \texttt{below\_floor} topics; no claim is made here about what
proportion of the full below-floor population falls into either category, only that both
categories are known to occur and are, in principle, distinguishable by inspection.

\section{Discussion: Implications for Combining OpenAlex with FOR/SEO-Coded
Administrative Data}
\label{sec:discussion}

For an analyst combining OpenAlex-derived publication data with administrative data
coded under FOR or SEO, this repository's contribution is threefold. First, it provides
a complete, versioned, auditable correspondence at whatever granularity the comparison
requires --- coarse comparisons can use the fully hand-curated
FOR2020-to-OpenAlex-field direction at confidence 1.0, while publication-level
classification can now use the OpenAlex-topic-to-FOR2020-field direction
(Section~\ref{sec:extension}), which reaches leaf-level precision for 87\% of topics and
degrades gracefully, rather than silently, for the remainder. Second, because every
correspondence carries an explicit confidence score and provenance tag, an analyst
chooses their own reliability threshold rather than inheriting an undocumented one: a
study requiring only division-level comparability can safely use confidence-0.0
group-level fallbacks unsuitable for a study requiring field-level precision. Third, the
residual, honestly-documented uncertainty in Section~\ref{sec:limitations} demonstrates
that automated cross-walks between independently-constructed classification systems
will always retain some genuinely ambiguous or unrecoverable cases, falling into at
least two qualitatively different categories --- fixable scoring errors versus
irreducible upstream classification mismatches --- with different implications for how
much further automated refinement can achieve.

The diagnostic approach used to identify these residual errors --- tracing a topic to
its subfield, that subfield to its matched FOR2020 group, and that group to its full
candidate-field list, then judging by inspection whether the correct answer is present
at all --- generalises to validating individual publication-level classifications in
downstream use. Because the topic-to-field correspondence is already pre-computed for
every OpenAlex topic, classifying an individual work requires no further
natural-language reasoning: its topic identifier is looked up directly against the
pre-built bridge table. Where an analyst wishes to go further --- sanity-checking a
low-confidence result against a specific work's own title, or attempting a better
classification for a work of particular importance --- the same short, structured
context (the work's title, optionally its abstract, and the small, pre-constrained list
of candidate fields from its topic's matched group) is directly suitable as input to
either a plausibility check (``is this classification consistent with this title?'') or
a re-classification judgement (``which candidate, if any, best fits this title?''),
including the judgement, illustrated by the biodiesel example above, that \emph{none} of
the candidates may be appropriate. Given corpus scale, such per-work verification is
realistically only practical as a targeted procedure applied to already-flagged
low-confidence cases; this repository's contribution is to make that minority
identifiable in the first place.

\section{Summary of Final Coverage}
\label{sec:results}

Table~\ref{tab:coverage} summarises the size of each canonical classification table and
the coverage achieved for the correspondences most relevant to combining OpenAlex with
FOR-coded administrative data, as of the work described in this document.

\begin{table}[htbp]
\centering
\small
\begin{tabular}{@{}lrrr@{}}
\toprule
\textbf{Correspondence} & \textbf{Source codes} & \textbf{Reaching finest level} & \textbf{Below-floor / division fallback} \\
\midrule
OpenAlex field $\rightarrow$ FOR2020 division & 26 & 26 (100\%) & 0 \\
OpenAlex subfield $\rightarrow$ FOR2020 group & 252 & 252 (100\%) & 0 \\
OpenAlex topic $\rightarrow$ FOR2020 field & 4{,}516 & 3{,}918 (86.8\%) & 598 (13.2\%, cascade to group) \\
FOR2020 $\rightarrow$ OpenAlex field/Leiden & 2{,}203 & 2{,}203 (100\%) & 0 \\
FOR1998 $\rightarrow$ OpenAlex field/Leiden & 1{,}060 & 1{,}060 (100\%) & 0 \\
FOR2008 $\rightarrow$ OpenAlex field/Leiden & 1{,}420 & 1{,}420 (100\%) & 0 \\
\bottomrule
\end{tabular}
\caption{Coverage of the OpenAlex$\leftrightarrow$FOR2020 correspondence and its
downstream legacy-vintage consequences, following the audit described in
Section~\ref{sec:extension}. ``Below-floor'' topic-level results never propagate to
\texttt{resolve()}'s output directly; they are recorded in the underlying bridge table
and the corresponding subfield-level (group) result is returned instead.}
\label{tab:coverage}
\end{table}

Across the repository as a whole, the canonical classification tables comprise 2{,}203
FOR2020 codes, 987 SEO2020 codes, 4{,}516 OpenAlex topics (nested under 252 subfields,
26 fields and 4 domains), 5 CWTS Leiden main fields, 361 ASJC codes and 22 SDG
goals/pillars, cross-linked by 13 bridge tables. The automated test suite comprises 29
tests, exercised against the pre-built database on every change, and passes in full as
of this document.

\section{Conclusion}

The extension described in this document closes what was, prior to this work, the most
significant remaining gap in the \texttt{ResearchClassification} repository's coverage:
the ability to classify an individual OpenAlex-indexed publication into a FOR2020 field
--- the level of granularity actually required to compare OpenAlex-derived research
output data against FOR-coded administrative collections meaningfully, rather than only
at the coarse disciplinary level the pre-existing correspondence supported. The
methodology used throughout --- direct, documented audit of algorithmically-generated
correspondences by inspection against their own candidate pools, in preference to
further automated formula-tuning, together with explicit, carried-through confidence
scoring and a documented rather than hidden residual error rate --- is offered here both
as a description of this specific repository's construction and as a template for
similar cross-walk exercises between other pairs of independently-constructed research
classification systems.

\end{document}
